\section*{Bab \ref{ch:b2:reaction-enthalpy} --- Entalpi Reaksi}

\begin{solution}{exo:b2:reaction-enthalpy:1}
$\Delta_r H^\circ = 3(-393.51) + 4(-285.83) - (-104.7) - 5(0) =
\qty{-2219.2}{kJ/mol}$, negatif: \omterm{def:b2:reaction-enthalpy:exothermic}{eksoterm}. (Entalpi pembakaran terukur
sama, \qty{-2219.2}{kJ/mol}.)
\end{solution}

\begin{solution}{exo:b2:reaction-enthalpy:2}
Metana: $890.3/16.04 = \qty{55.5}{kJ/g}$, jadi \qty{55.5}{MJ} per
kilogram. Dihidrogen: $285.83/2.016 = \qty{141.8}{kJ/g}$,
\qty{141.8}{MJ/kg}: sekitar 2.6 kali lebih banyak per kilogram (tetapi
jauh lebih sedikit per liter, karena gasnya ringan).
\end{solution}

\begin{solution}{exo:b2:reaction-enthalpy:3}
$\Delta\nu_{\text{gas}} = 6 - 7.5 = -1.5$, sehingga $\Delta_r H = \Delta_r U +
\Delta\nu_{\text{gas}}RT = -3264 - 1.5 \times 8.314 \times 298.15 \times
10^{-3} = \qty{-3267.7}{kJ/mol}$, sesuai dengan nilai yang dihitung dari
entalpi pembentukan, \qty{-3267.6}{kJ/mol}.
\end{solution}

\begin{solution}{exo:b2:reaction-enthalpy:4}
Reaksi sasaran adalah reaksi pertama dikurangi dua kali reaksi kedua:
$\Delta_r H^\circ = -393.5 - 2(-283.0) = \qty{+172.5}{kJ/mol}$. Reaksi ini
\omterm{def:b2:reaction-enthalpy:exothermic}{endoterm}: kokas panas menjadi dingin ketika mengubah karbon dioksida
menjadi karbon monoksida.
\end{solution}

\begin{solution}{exo:b2:reaction-enthalpy:5}
Ikatan yang diputus: C--H dan Cl--Cl, $416 + 243 = \qty{659}{kJ/mol}$;
ikatan yang terbentuk: C--Cl dan H--Cl, $350 + 432 = \qty{782}{kJ/mol}$:
perkiraannya \qty{-123}{kJ/mol}. Dari entalpi pembentukan:
$-83.7 - 92.3 + 74.9 =
\qty{-101.1}{kJ/mol}$. Selisih \qty{22}{kJ/mol}
berasal dari nilai rata-rata C--Cl, yang dirata-ratakan atas beberapa
molekul, dan dari ikatan C--H yang diputus pada metana, yang lebih kuat
daripada rata-rata keempatnya.
\end{solution}

\begin{solution}{exo:b2:reaction-enthalpy:6}
Pada \qty{298}{K}: $-241.83 + 285.83 = \qty{44.00}{kJ/mol}$. Kirchhoff
dengan $\Delta C_p = 33.59 - 75.35 = \qty{-41.76}{J/(K.mol)}$:
$44.00 - 0.04176 \times 101.85 = \qty{39.75}{kJ/mol}$ pada \qty{400}{K}.
Tabel memberikan $-242.85 + 282.59 = \qty{39.74}{kJ/mol}$: kapasitas kalor
tetap sangat baik dalam selang ini.
\end{solution}

\begin{solution}{exo:b2:reaction-enthalpy:7}
Ionisasi: $4.341 \times 96.485 = \qty{418.8}{kJ/mol}$; penangkapan
elektron: $-3.613 \times 96.485 = \qty{-348.6}{kJ/mol}$. \omterm{def:b2:reaction-enthalpy:lattice-enthalpy}{Entalpi kisi}
$= 436.7 +
89.0 + 121.3 + 418.8 - 348.6 = \qty{717}{kJ/mol}$, lebih kecil
daripada \qty{787}{kJ/mol} untuk natrium klorida: \ce{K+} lebih besar
daripada \ce{Na+}, ion-ionnya berjauhan dan saling menarik lebih lemah.
\end{solution}

\begin{solution}{exo:b2:reaction-enthalpy:8}
Kapasitas kalor kalorimeter: $C = 2.00/0.418 = \qty{4.785}{kJ/K}$. Kalor
yang dilepaskan: $4.785 \times 0.583 = \qty{2.789}{kJ}$, untuk
$\xi =
\qty{0.0500}{mol}$ (asamnya pembatas, basanya sedikit berlebih):
$\Delta_r H = -2.789/0.0500 = \qty{-55.8}{kJ/mol}$.
\end{solution}

\begin{solution}{exo:b2:reaction-enthalpy:9}
\ce{C8H18(l) + 25/2O2(g) -> 8CO2(g) + 9H2O(l)}: $\Delta_r H^\circ =
8(-393.51) + 9(-285.83) + 250.3 = \qty{-5470.3}{kJ/mol}$; dengan
$M =
\qty{114.2}{g/mol}$, \qty{47.9}{MJ/kg}. Karbon dioksida:
$8 \times 44.01 =
\qty{352.1}{g}$ untuk \qty{5.470}{MJ}, yaitu
\qty{64.4}{g/MJ}, dibandingkan dengan $44.01/0.8903 = \qty{49.4}{g/MJ}$
untuk metana: metana mempunyai lebih banyak hidrogen per karbon, dan
pembakaran hidrogen melepaskan kalor tanpa karbon dioksida.
\end{solution}

\begin{solution}{exo:b2:reaction-enthalpy:10}
Produk \ce{CH4 + 2O2 -> CO2 + 2H2O(g)}: $C_p = 37.13 + 2 \times 33.59 =
\qty{104.3}{J/K}$; $q = \qty{802.3}{kJ}$; $T_f = 298 + 802\,300/104.3 \approx
\qty{7990}{K}$, hampir tiga kali nilai di udara, karena tidak ada nitrogen
yang mengambil bagian kalornya. Nyala oksigen yang sesungguhnya jauh
lebih dingin: jauh sebelum \qty{7990}{K}, karbon dioksida dan air
terdisosiasi menjadi karbon monoksida, hidrogen, oksigen, dan atom-atom,
reaksi-reaksi \omterm{def:b2:reaction-enthalpy:exothermic}{endoterm} yang menyerap sebagian besar kalor; selain itu,
kapasitas kalor bertambah seiring $T$.
\end{solution}

\begin{solution}{exo:b2:reaction-enthalpy:11}
$\Delta_r H^\circ(700) = \Delta_r H^\circ(298) + \int_{298}^{700}(\Delta_r a +
\Delta_r b\,T)\,\dd T = -91.88 + [-60.5 \times 401.85 + 0.0270 \times
(700^2 - 298.15^2)] \times 10^{-3} = -91.88 - 24.31 + 10.83 =
\qty{-105.4}{kJ/mol}$, hanya berbeda \qty{0.2}{kJ/mol} dari tabel
(\qty{-105.2}{kJ/mol}), sedangkan perkiraan $C_p$ tetap meleset
\qty{4.5}{kJ/mol}.
\end{solution}

\begin{solution}{exo:b2:reaction-enthalpy:12}
$\Delta_r H^\circ = 2 \times 90.29 = \qty{+180.6}{kJ/mol}$ (\omterm{def:b2:reaction-enthalpy:exothermic}{endoterm}).
$\Delta_r C_p^\circ = 2(29.85) - 29.12 - 29.38 = \qty{1.2}{J/(K.mol)}$,
sehingga antara \qty{298}{K} dan \qty{2000}{K} \omterm{def:b2:reaction-enthalpy:reaction-quantity}{entalpi reaksi} berubah
sekitar $1.2 \times 1702 = \qty{2}{kJ/mol}$, satu persen. Reaksi \omterm{def:b2:reaction-enthalpy:exothermic}{endoterm}
diuntungkan oleh suhu tinggi (\cref{ch:b2:equilibrium-shifts}): udara
dingin hampir tidak menghasilkan nitrogen monoksida, nyala yang paling
panas menghasilkan paling banyak; itulah sebabnya suhu pembakaran menjadi
pengungkit untuk menekan polutan ini.
\end{solution}

\begin{solution}{pb:b2:reaction-enthalpy:1}
\textbf{1.} \ce{C4H10(g) + 13/2O2(g) -> 4CO2(g) + 5H2O(l)}.
\textbf{2.} $\Delta_r H^\circ = 4(-393.51) + 5(-285.83) + 125.6 =
\qty{-2877.6}{kJ/mol}$.
\textbf{3.} Dengan uap air, $4(-393.51) + 5(-241.83) + 125.6 =
\qty{-2657.6}{kJ/mol}$; selisihnya, \qty{220.0}{kJ}, adalah entalpi
penguapan lima mol air (masing-masing \qty{44.0}{kJ/mol}).
\textbf{4.} $2877.6/58.12 = \qty{49.5}{kJ/g}$ dan $2657.6/58.12 =
\qty{45.7}{kJ/g}$.
\textbf{5.} Nilai kedua: air keluar sebagai uap dan kalor pengembunannya
hilang.
\textbf{6.} $230 \times 45.7 = \qty{1.05e4}{kJ}$, sekitar \qty{10.5}{MJ}.
\textbf{7.} Pada volume tetap, kalor yang diterima adalah $\Delta U$ (tidak
ada kerja tekanan): bom mengukur $\Delta_r U$.
\textbf{8.} $\Delta\nu_{\text{gas}} = 4 - 1 - 6.5 = -3.5$ (airnya berupa
cairan).
\textbf{9.} $\Delta_r U = \Delta_r H - \Delta\nu_{\text{gas}}RT = -2877.6 +
3.5 \times 2.479 = \qty{-2868.9}{kJ/mol}$.
\textbf{10.} $\xi = 0.500/58.12 = \qty{8.60e-3}{mol}$; kalor $8.60 \times
10^{-3} \times 2868.9 = \qty{24.68}{kJ}$; $\Delta T = 24.68/10.00 =
\qty{2.468}{K}$.
\textbf{11.} Kalor $10.00 \times 2.461 = \qty{24.61}{kJ}$: $\Delta_r U =
-24.61/(8.603 \times 10^{-3}) = \qty{-2860.6}{kJ/mol}$ dan $\Delta_r H =
-2860.6 - 8.7 = \qty{-2869.3}{kJ/mol}$, \qty{0.29}{\percent} di bawah nilai
tabel dalam nilai mutlak: kalor yang hilang dari kalorimeter atau
pembakaran yang sedikit tidak sempurna.
\textbf{12.} $n = 1000/18.02 = \qty{55.5}{mol}$; $Q = 55.5 \times 75.35
\times 80 = \qty{334.5}{kJ}$.
\textbf{13.} Kalor yang dilepaskan $16.0 \times 45.72 = \qty{731.5}{kJ}$;
efisiensi $334.5/731.5 = 0.46$.
\textbf{14.} Ke gas-gas pembakaran panas yang lolos di sekeliling panci,
udara yang dipanaskan oleh nyala, dinding panci, dan radiasi.
\textbf{15.} $230 \times 45.72 \times 0.457/334.5 = \qty{14.4}{L}$.
\textbf{16.} \qty{6.5}{mol} \ce{O2} disertai $6.5 \times 78.08/20.95 =
\qty{24.23}{mol}$ \ce{N2} dan $6.5 \times 0.93/20.95 = \qty{0.289}{mol}$
\ce{Ar}.
\textbf{17.} \qty{4}{mol} \ce{CO2}, \qty{5}{mol} \ce{H2O(g)},
\qty{24.23}{mol} \ce{N2}, \qty{0.289}{mol} \ce{Ar}.
\textbf{18.} Nyala itu adiabatik dan isobarik, sehingga $\Delta H = 0$.
Sepanjang jalur yang terdiri atas reaksi pada \qty{298}{K}
($\Delta H_1 =
\qty{-2657.6}{kJ}$) diikuti pemanasan produk-produk ini,
$\Delta H_2 =
-\Delta H_1 = \qty{+2657.6}{kJ}$: seluruh kalor masuk ke
gas-gas itu.
\textbf{19.} $C_p = 4(37.13) + 5(33.59) + 24.23(29.12) + 0.289(20.79) =
\qty{1028}{J/K}$.
\textbf{20.} $T_f = 298 + 2\,657\,600/1028 = \qty{2883}{K}$.
\textbf{21.} $T_f = 2300 + 100 \times (2657.6 - 2515.7)/(2655.7 - 2515.7) =
\qty{2401}{K}$.
\textbf{22.} Kapasitas kalor gas bertambah seiring suhu (sekitar
seperempat antara \qty{298}{K} dan \qty{2400}{K}): nilai pada suhu kamar
memanaskan gas-gas itu terlalu murah.
\textbf{23.} Di atas \qty{2000}{K}, karbon dioksida dan air sebagian
terdisosiasi (\omterm{def:b2:reaction-enthalpy:exothermic}{endoterm}), nyala memancarkan radiasi, dan pencampuran dengan
udara tidak pernah sempurna.
\textbf{24.} \omterm{def:b2:reaction-enthalpy:flame-temperature}{Suhu nyala adiabatik} butana di udara adalah
$\boldsymbol{T_{\text{ad}} \approx \qty{2.40e3}{K}}$.
\end{solution}
